\documentclass[11pt]{article}
\usepackage[margin=1in]{geometry}
\usepackage{amsmath,amssymb,amsthm}
\newtheorem{theorem}{Theorem}
\newtheorem{lemma}[theorem]{Lemma}
\newtheorem{proposition}[theorem]{Proposition}
\newtheorem{corollary}[theorem]{Corollary}
\theoremstyle{definition}
\newtheorem{remark}[theorem]{Remark}

\title{Structural constraints on covering systems with all moduli of the form $p-1$,\\
$p\ge5$ prime\\[2mm]
\large (Erd\H{o}s Problem 273: partial results; the problem is \emph{not} resolved here)}
\author{}
\date{}

\begin{document}
\maketitle

\begin{abstract}
Erd\H{o}s and Graham asked whether there is a covering system of $\mathbb Z$ all of whose moduli
are of the form $p-1$ for primes $p\ge5$. We do \emph{not} answer this. We record a set of
rigorously proved structural constraints and reductions, each independently audited and
machine-verified, together with an explicit determination of why several natural attacks cannot
work. In particular we prove: an exact parity-split equivalence reducing the problem to the
existence of two \emph{disjoint} covering systems inside the halved set $H=\{m\ge2:2m+1\text{ prime}\}$;
that the infimum of $\sum 1/n_i$ over \emph{unrestricted} distinct-moduli covering systems is
exactly $1$, so that no universal constant $c>1$ with $\sum 1/n_i\ge c$ exists to be exploited;
a rigidity theorem identifying the dyadic staircase as the unique exact cover of $\mathbb Z$ minus
one class by distinct moduli, which no subset of $E$ can realise; and two effective elimination
tests (a coprimality/pigeonhole test and an exact $q$-adic fiber test) which together rule out
explicit lists of possible least common multiples. All claims are accompanied by independent
verification scripts.
\end{abstract}

\section{Setting}

A \emph{covering system} is a finite list of residue classes $a_i \pmod{n_i}$, $1\le i\le k$, with
$1<n_1<n_2<\cdots<n_k$ (so the moduli are pairwise distinct and exceed $1$) and
$\bigcup_i (a_i+n_i\mathbb Z)=\mathbb Z$. Put
\[
E:=\{p-1:\ p\ \text{prime},\ p\ge5\}=\{n\ge4:\ n+1\ \text{prime}\},\qquad
H:=\{m\ge2:\ 2m+1\ \text{prime}\}=\tfrac12E .
\]
Every $n\in E$ is even, and $n\mapsto n/2$ is a bijection $E\to H$. We write
$B_E(X)=\sum_{n\in E,\ n\le X}1/n$ and, for $L\ge1$, $\beta(L)=\sum_{n\mid L,\ n\in E}1/n$.
(We reserve the letter $f$ for the functional of Lemma~\ref{lem:A2}.)

\smallskip
\noindent\textbf{Question (Erd\H{o}s--Graham, 1980).} Is there a covering system with all moduli in
$E$?

\section{The parity split}

\begin{lemma}[Parity split]\label{lem:parity}
The following are equivalent.
\begin{enumerate}
\item[\rm(i)] There is a covering system with all moduli in $E$.
\item[\rm(ii)] There are two \emph{disjoint} finite nonempty sets $M_0,M_1\subseteq H$ such that for
each $j$ there is a covering system of $\mathbb Z$ whose set of moduli is exactly $M_j$.
\end{enumerate}
Moreover $\operatorname{lcm}(E\text{-moduli})=2\operatorname{lcm}(H\text{-moduli})$.
\end{lemma}

\begin{proof}
(i)$\Rightarrow$(ii). Let $\{a_i\bmod n_i\}$ be such a system, $n_i=2m_i$ with $m_i\in H$; the
$m_i$ are pairwise distinct because the $n_i$ are. Since $n_i$ is even, $a_i+n_i\mathbb Z$ consists
solely of integers of the parity of $a_i$; so $I_j=\{i:a_i\equiv j\ (2)\}$ partitions the index
set, and $M_j=\{m_i:i\in I_j\}$ are disjoint. Fix $j$ and $y\in\mathbb Z$; put $x=2y+j$. Some $i$
has $x\equiv a_i\ (2m_i)$, forcing $a_i\equiv j\ (2)$, i.e. $i\in I_j$; writing $a_i=2b_i+j$,
\[
2y+j\equiv 2b_i+j \pmod{2m_i}\iff 2(y-b_i)\in 2m_i\mathbb Z\iff y\equiv b_i\pmod{m_i},
\]
so $\{b_i\bmod m_i\}_{i\in I_j}$ covers $\mathbb Z$. (The argument quantifies over all
$y\in\mathbb Z$, negative ones included; classes are two-sided.) Each $M_j$ is nonempty, since
otherwise no class would cover the integers of parity $j$.

(ii)$\Rightarrow$(i). Given coverings $\{b^{(j)}_m\bmod m\}_{m\in M_j}$, take the classes
$2b^{(j)}_m+j \pmod{2m}$ for all $m\in M_j$, $j=0,1$. The moduli $2m$ lie in $E$ and are pairwise
distinct: within a block because the $m$ are, across blocks because $M_0\cap M_1=\varnothing$.
(The disjointness hypothesis is essential and not decorative: $M_0=M_1=\{2\}$ would produce the
modulus list $4,4$, violating distinctness.) If $x=2y+j$ then $y\equiv b^{(j)}_m\ (m)$ for some
$m\in M_j$, whence $x\equiv 2b^{(j)}_m+j\ (2m)$.
\end{proof}

\begin{corollary}\label{cor:pivot}
Since $2,3\in H$ each occur once, at most one of $M_0,M_1$ contains $2$ and at most one contains
$3$. Hence an affirmative answer requires, in particular, a covering system of $\mathbb Z$ with
distinct moduli all of the form $(p-1)/2$, $p\ge7$ prime, i.e.\ an $H$-covering avoiding the
modulus $2$. Contrapositively: if no such system exists, the answer to the problem is negative.
\end{corollary}

\begin{remark}
Any construction forcing one of $M_0,M_1$ to consist of odd moduli only is a dead end: that is the
Erd\H{o}s--Selfridge odd covering problem, which is open.
\end{remark}

\section{The infimum of the reciprocal sum}

\begin{lemma}[Doubling map]\label{lem:double}
If $C=\{a_i\bmod n_i\}$ is a covering system with distinct moduli, so is
$D(C)=\{0\bmod2\}\cup\{(2a_i+1)\bmod 2n_i\}$, and
$\operatorname{cost}(D(C))=\tfrac12+\tfrac12\operatorname{cost}(C)$, where
$\operatorname{cost}=\sum_i 1/n_i$.
\end{lemma}

\begin{proof}
Even $x$ lie in $0\bmod 2$. If $x=2y+1$, pick $i$ with $y\equiv a_i\ (n_i)$; then
$x\equiv 2a_i+1\ (2n_i)$. The moduli are $2$ and the pairwise distinct $2n_i\ge4$.
\end{proof}

\begin{theorem}\label{thm:inf}
$\inf\sum_i 1/n_i=1$ over covering systems with distinct moduli, and the infimum is not attained.
Explicitly, iterating Lemma~\ref{lem:double} from $C_0=\{0(2),0(3),1(4),5(6),7(12)\}$ gives systems
$C_k$ with $\sum 1/n_i=1+\tfrac13 2^{-k}$.
\end{theorem}

\begin{proof}
Subadditivity of density gives $\sum 1/n_i\ge1$; equality would make the system an exact cover
with distinct moduli (two classes are disjoint or meet in a class of positive density), which is
impossible by Davenport--Mirsky--Newman--Rado. The upper bound is Lemma~\ref{lem:double}.
\end{proof}

\begin{remark}[what this does and does not give]\label{rem:budget}
Theorem~\ref{thm:inf} concerns \emph{unrestricted} distinct-moduli systems. It shows that no
universal constant $c>1$ with $\sum_i 1/n_i\ge c$ exists, so any attack that would derive a
contradiction from an absolute lower bound on the reciprocal sum is dead. It does \emph{not} say
that $E$- or $H$-systems can have cost close to $1$: no lower bound on the cost of an $E$- or
$H$-covering is known, and the cheapest $H$-covering found in extensive search costs $65/48$ with
no downward trend. Indeed $D$ never produces an $E$-system: every $D(C)$ contains the modulus
$2\notin E$, and admissibility of $n$ and of $2n$ are unrelated conditions.
\end{remark}

\begin{remark}[a dichotomy, on unverified authority]\label{rem:FK}
Filaseta and Kalogirou (arXiv:2407.15280, ``Covering systems with the sum of the reciprocals of the
moduli close to $1$'') are reported to prove the 1973 Erd\H{o}s--Selfridge conjecture that a least
modulus exceeding $4$ forces $\sum 1/n_i$ to be bounded away from $1$. \emph{We were unable to
open the primary source} (network egress was blocked); the statement was confirmed only at
abstract level and must not be treated as established. If it is as reported, it yields a
dichotomy rather than a blanket conclusion: an $E$-covering that \emph{uses} the modulus $4$ has
least modulus $4$ and lies in the regime where no excess bound is available, whereas an
$E$-covering \emph{avoiding} $4$ has least modulus $\ge 6$ and would inherit a positive excess
bound. We do not use either horn below.
\end{remark}

\section{Rigidity of the cheap mechanism}

\begin{theorem}[Rigidity]\label{thm:rigid}
If residue classes $a_i\bmod n_i$ ($0\le a_i<n_i$) with pairwise distinct moduli $>1$ partition
$\mathbb Z\setminus(c\bmod D)$ exactly, then $D=2^m$ and the moduli are exactly $2,4,\dots,2^m$.
\end{theorem}

\begin{proof}[Proof]
For $|z|<1$, $\sum_i z^{a_i}/(1-z^{n_i})=\frac1{1-z}-\frac{z^{c}}{1-z^{D}}$. Let $n_k$ be the
largest modulus. At a primitive $n_k$-th root of unity the left side has a pole (only $i=k$ can
contribute), so the right side must too, giving $n_k\mid D$; at a primitive $D$-th root the right
side has a pole, so some $n_i$ is divisible by $D$, giving $D\mid n_i\le n_k$. Hence $n_k=D$.
Matching residues at the primitive $D$-th roots forces $\zeta^{a_k-c}=-1$ for all of them, so $D$
is even and $a_k\equiv c+D/2$. Deleting that class leaves an exact partition of
$\mathbb Z\setminus(c\bmod D/2)$ by the remaining classes; induct, with base $D=2$.
\end{proof}

\begin{corollary}\label{cor:rigid}
$\{4,16,256,65536\}\subseteq E\cap\{2^k\}$ and $\{2,8,128,32768\}\subseteq H\cap\{2^k\}$; these are
the known members, coming from the known Fermat primes. Unconditionally, however,
$2\notin E$ and $4\notin H$ (since $9=3^2$), and $E$ contains no two consecutive powers of $2$:
$2^k+1$ prime forces $k$ to be a power of $2$, and $k,k+1$ are both powers of $2$ only for $k=1$,
which gives $2\notin E$. Consequently \emph{no} non-empty dyadic staircase lies in $E$, and the
only one inside $H$ is the trivial $\{2\}$. So the mechanism of Theorem~\ref{thm:rigid} is
unavailable in $E$ and essentially unavailable in $H$.
\end{corollary}

\begin{remark}
Corollary~\ref{cor:rigid} removes the \emph{only known} mechanism for driving the reciprocal cost
to $1$; it does not exclude other mechanisms, and in particular it is not a lower bound on the cost
of an $E$- or $H$-covering. No such lower bound is proved anywhere in this note.
\end{remark}

\section{Two effective elimination tests}

\begin{lemma}[Forced overlap]\label{lem:A2}
Let $M$ be the modulus set of a covering system with distinct moduli and let $T\subseteq M$ be
pairwise coprime. Then
\[
\sum_{m\in M}\frac1m-1\ \ge\ f(T):=\sum_{m\in T}\frac1m-1+\prod_{m\in T}\Bigl(1-\frac1m\Bigr).
\]
Both hypotheses are needed: $T\subseteq M$ and pairwise coprimality each fail the conclusion
otherwise.
\end{lemma}

\begin{proof}
For a subfamily $S$ of the chosen classes put $g(S)=\sum_{m\in S}1/m-\operatorname{dens}(\bigcup S)$;
the density exists because a finite union of classes is periodic modulo the lcm. Adding a class of
a new modulus $m'$ raises the first term by $1/m'$ and the density by at most $1/m'$, so $g$ is
non-decreasing. $g(M)=\sum_M 1/m-1$ since the union is $\mathbb Z$. For pairwise coprime $T$ the
classes are independent by the Chinese remainder theorem, \emph{whatever the residues}: the
complement has exactly $\prod_{m\in T}(m-1)$ residues modulo $\prod_{m\in T}m$, so
$\operatorname{dens}(\bigcup T)=1-\prod_{m\in T}(1-1/m)$ and $g(T)=f(T)$.
\end{proof}

\begin{theorem}\label{thm:A3}
If $60\mid L$ and $1<\beta(L)\le\frac{31}{30}$, then no covering system with all moduli in $E$ has
lcm dividing $L$.
\end{theorem}

\begin{proof}
$60\mid L$ gives $4,6,10\mid L$, all in $E$. Since $\beta(L)-1\le\frac1{30}<\frac1{10}$, each of
$4,6,10$ is forced to be used: deleting it leaves reciprocal sum $<1$. Their halves $2,3,5$ then
lie in $M_0\cup M_1$, which is what licenses $T\subseteq M_j$ in Lemma~\ref{lem:A2}. With
$X_j=\sum_{m\in M_j}1/m-1$, disjointness gives $X_0+X_1=\sum_{M_0\cup M_1}1/m-2\le 2\beta(L)-2
\le\frac1{15}$, while both $X_j>0$ strictly (Theorem~\ref{thm:inf}); hence $X_j<\frac1{15}$ for
each $j$. As $f(\{2,3\})=\frac16$, $f(\{2,5\})=\frac1{10}$, $f(\{3,5\})=\frac1{15}$,
Lemma~\ref{lem:A2} forbids any two of $2,3,5$ from lying in the same $M_j$. Three objects cannot
occupy two blocks.
\end{proof}

\begin{remark}
The inequality is tight: $f(\{3,5\})$ \emph{equals} the bound on $X_0+X_1$, so the proof survives
only because both strict inequalities above are strict. Theorem~\ref{thm:A3} is a statement about
a fixed $L$, never about all finite subsets of $E$.
\end{remark}

\begin{theorem}[strengthened form]\label{thm:A3plus}
Fix $L$ even, put $S=\{m\mid L/2:\ m\in H\}$ and $B=\sum_{m\in S}1/m$. Call $m\in S$
\emph{forced} if $B-1/m<2$, and let $F$ be the set of forced moduli. For $S'\subseteq S$ let
$g^*(S')=\max\{f(T):T\subseteq S'\text{ pairwise coprime}\}$. If
\[
\min_{F=F_0\sqcup F_1}\bigl(g^*(F_0)+g^*(F_1)\bigr)\ >\ B-2 ,
\]
then no covering system with all moduli in $E$ has lcm $L$.
\end{theorem}

\begin{proof}
Every used modulus lies in $S$; if $m$ were unused the two halves would have total reciprocal mass
at most $B-1/m<2$, contradicting $X_0,X_1>0$. So $F\subseteq M_0\cup M_1$ and the halves induce a
2-colouring of $F$. By Lemma~\ref{lem:A2}, $X_j\ge g^*(F_j)$, while $X_0+X_1\le B-2$.
\end{proof}

\noindent Theorem~\ref{thm:A3} is the case $F\supseteq\{2,3,5\}$ with $|T|=2$.

\begin{lemma}[$q$-adic fiber test]\label{lem:fiber}
Let $q$ be a prime, and let $S$ be a set of $H$-moduli (so every element is $\ge2$). For an
assignment of residues $b_m$ and $r\in\mathbb Z_q$ put
$F_q(r)=\sum_{m:\,b_m\equiv r\ (q^{\nu_q(m)})} q^{\nu_q(m)}/m$. If some covering system with all
moduli in $E$ has $\operatorname{lcm}/2=L_H$ and $S=\{m\mid L_H:\ m\ge2,\ 2m+1\text{ prime}\}$,
then $\Phi_q(S):=\max_{\text{assignments}}\min_r F_q(r)\ \ge\ 2$ for every prime $q\mid L_H$.
\end{lemma}

\begin{proof}
For a single covering system drawn from $S$, and $J\ge\max_m\nu_q(m)$: since
$\gcd(m,q^J)=q^{\nu_q(m)}$, the integers $y\equiv r\ (q^J)$ that the class $b_m\ (m)$ meets are
exactly those with $b_m\equiv r\ (q^{\nu_q(m)})$, and inside that fiber the class has relative
density $q^{\nu_q(m)}/m$; subadditivity forces $F_q(r)\ge1$. By Lemma~\ref{lem:parity} the two
halves are disjoint subsets of $S$ carrying their own residues, so their fiber functions add.
Enlarging the used set to all of $S$ only increases $F_q$ pointwise.
\end{proof}

\begin{remark}
With $\mu$ the Haar measure on $\mathbb Z_q$, $\int F_q\,d\mu=\sum_{m\in S}1/m$ exactly, the sum
being over the \emph{$H$}-moduli. So the reciprocal-sum condition is only the \emph{average} of the
fiber condition, whereas covering demands the \emph{minimum}.
\end{remark}

\section{What is ruled out, and what is not}

Combining $\beta(L)>1$, Theorem~\ref{thm:A3}, Theorem~\ref{thm:A3plus} and Lemma~\ref{lem:fiber}
(the last computed exactly by branch and bound) eliminates explicit candidate least common
multiples; the computations and their independent verifications are in the accompanying scripts.
The smallest $L$ with $\beta(L)>1$ is $55440=2^4\cdot3^2\cdot5\cdot7\cdot11$, and $55440$ is itself
eliminated twice independently: by Lemma~\ref{lem:fiber} at $q=2$, and by
Theorem~\ref{thm:A3plus} (the forced set is $\{2,3,5,6,8,9,11\}$ and every 2-colouring gives
$g^*(F_0)+g^*(F_1)\ge 6/55>269/3080=B-2$, a witness colouring being
$\{3,5,6,9,11\}\mid\{2,8\}$ with $T=\{3,5,11\}$).

Finally we record an obstruction to the obstructions. Fix a finite set $Q$ of primes and consider
any elimination test built from monotone additive fiber budgets over the cells of
$\prod_{q\in Q}\mathbb Z_q$. Any $m$ coprime to $\prod Q$ has $\nu_q(m)=0$ for all $q\in Q$ and so
contributes $1/m$ to every cell for every residue assignment; and
$\sum_{m\in H,\ \gcd(m,\prod Q)=1}1/m=\infty$, being $\sum 2/(p-1)$ over the primes $p$ in a union
of admissible residue classes (Mertens--Dirichlet). Hence for each fixed $Q$ there are arbitrarily
large pools passing the $Q$-local test trivially, so no \emph{fixed}-$Q$ test can eliminate all
lattices. (The eliminations above are not refuted by this: they use a prime $q$ depending on $L$.)
A negative resolution therefore appears to need a functional coupling all primes with a bound
uniform in their number. The available technology of that kind --- Hough's and
Balister--Bollob\'as--Morris--Sahasrabudhe--Tiba's distortion method --- proves that the least
modulus of a distinct-moduli covering system is at most $10^{16}$, respectively $616000$; it
yields a contradiction only for systems all of whose moduli exceed such a bound, whereas
$\min E=4$. So it does not apply here.

\section*{Status}

The problem is \textbf{not resolved} by this note. The results above are constraints and
reductions only. The decisive open question is Corollary~\ref{cor:pivot}: whether $\mathbb Z$
admits a covering system with distinct moduli all in $H\setminus\{2\}$.

\end{document}
